\begin{frame}{Definición inductiva del determinante}

Sea \(A=(a_{ij})\in M_n(k)\). Denotaremos por \(A_{ij}\) la matriz
de orden \(n-1\) que se obtiene de \(A\) suprimiendo la fila \(i\)
y la columna \(j\).

\medskip

\begin{center}
\begin{tikzpicture}[baseline=(M.center)]
\matrix (M) [
    matrix of math nodes,
    ampersand replacement=\&,
    left delimiter=(,
    right delimiter=),
    row sep=2mm,
    column sep=3mm
] {
a_{11} \& \cdots \& a_{1j} \& \cdots \& a_{1n}\\
\vdots \&        \& \vdots \&        \& \vdots\\
a_{i1} \& \cdots \& a_{ij} \& \cdots \& a_{in}\\
\vdots \&        \& \vdots \&        \& \vdots\\
a_{n1} \& \cdots \& a_{nj} \& \cdots \& a_{nn}\\
};

\draw[ultra thick]
    ([xshift=-2mm]M-3-1.west)
    --
    ([xshift=2mm]M-3-5.east);

\draw[ultra thick]
    ([yshift=2mm]M-1-3.north)
    --
    ([yshift=-2mm]M-5-3.south);

\node[right=7mm of M] {\(= A_{ij}\)};
\end{tikzpicture}
\end{center}
\end{frame}

\begin{frame}{Definición inductiva del determinante}

\begin{definition}
Definimos inductivamente el \alert{determinante} de una matriz
cuadrada.

Si \(A=(a)\in M_1(k)\), definimos
\[
\det(A)=a.
\]

Si \(n>1\), suponiendo definido el determinante para matrices de
orden \(n-1\), para \(A=(a_{ij})\in M_n(k)\) definimos
\[
\det(A)
=
\sum_{j=1}^n(-1)^{1+j}a_{1j}\det(A_{1j}).
\]
Esta fórmula se conoce como el \alert{desarrollo del determinante por la primera fila}.

El determinante de una matriz $A$ también se denotará por 
\[|A|=\det(A).\]
\end{definition}

\end{frame}

\begin{frame}{Determinantes de orden \(2\)}

Sea
\[
A=
\begin{pmatrix}
a&b\\
c&d
\end{pmatrix}.
\]

Por la definición inductiva,
\[
\begin{aligned}
\det(A)
&=
a\det(d)-b\det(c)\\
&=ad-bc.
\end{aligned}
\]

Por tanto,
\[
\det
\begin{pmatrix}
a&b\\
c&d
\end{pmatrix}
=
ad-bc.
\]

\end{frame}

\begin{frame}[allowframebreaks]{Determinantes de orden \(3\)}

Sea
\[
A=
\begin{pmatrix}
a_{11}&a_{12}&a_{13}\\
a_{21}&a_{22}&a_{23}\\
a_{31}&a_{32}&a_{33}
\end{pmatrix}.
\]

Desarrollando por la primera fila,
\[
\begin{aligned}
\det(A)
={}&
a_{11}
\begin{vmatrix}
a_{22}&a_{23}\\
a_{32}&a_{33}
\end{vmatrix}
-
a_{12}
\begin{vmatrix}
a_{21}&a_{23}\\
a_{31}&a_{33}
\end{vmatrix}\\
&+
a_{13}
\begin{vmatrix}
a_{21}&a_{22}\\
a_{31}&a_{32}
\end{vmatrix}.
\end{aligned}
\]

Usando la fórmula para determinantes de orden \(2\),
\[
\begin{aligned}
\det(A)
={}&
a_{11}a_{22}a_{33}
-a_{11}a_{23}a_{32}\\
&-a_{12}a_{21}a_{33}
+a_{12}a_{23}a_{31}\\
&+a_{13}a_{21}a_{32}
-a_{13}a_{22}a_{31}.
\end{aligned}
\]
\end{frame}

\begin{frame}[allowframebreaks]{Determinantes de orden \(3\)}

\begin{columns}
\begin{column}{0.5\textwidth}
Reordenando los términos obtenemos la \alert{regla de Sarrus}:
\[
\begin{aligned}
\det(A)
={}&
a_{11}a_{22}a_{33}
+a_{12}a_{23}a_{31}
+a_{13}a_{21}a_{32}\\
&-
a_{13}a_{22}a_{31}
-a_{11}a_{23}a_{32}
-a_{12}a_{21}a_{33}.
\end{aligned}
\]
Esta regla solo es válida para matrices de orden \(3\).
\end{column}

\begin{column}{0.5\textwidth}
% Source - https://tex.stackexchange.com/a/603570
% Posted by SebGlav
% Retrieved 2026-09-24, License - CC BY-SA 4.0
\begin{center}
\begin{tikzpicture}[baseline=(A.center)]
  \tikzset{BarreStyle/.style =   {opacity=.4,line width=4 mm,line cap=round,color=#1}}
    \tikzset{SignePlus/.style =   {above left,inner sep=1.5pt,opacity=.4,circle,fill=#1}}
    \tikzset{SigneMoins/.style =   {below left,inner sep=1.5pt,opacity=.4,circle,fill=#1}}
% the matrices
\matrix (A) [matrix of math nodes, nodes = {node style ge},column sep=0 mm, ampersand replacement=\&,] 
{ a_{11} \& a_{12} \& a_{13}  \\
  a_{21} \& a_{22} \& a_{23}  \\
  a_{31} \& a_{32} \& a_{33}  \\
  a_{11} \& a_{12} \& a_{13} \\
  a_{21} \& a_{22} \& a_{13}\\
};

 \draw [BarreStyle=blue] (A-1-1.north west) node[SignePlus=blue] {$+$} to (A-3-3.south east) ;
 \draw [BarreStyle=blue] (A-2-1.north west) node[SignePlus=blue] {$+$} to (A-4-3.south east) ;
 \draw [BarreStyle=blue] (A-3-1.north west) node[SignePlus=blue] {$+$} to (A-5-3.south east) ;
 \draw [BarreStyle=red]  (A-3-1.south west) node[SigneMoins=red] {$-$} to (A-1-3.north east);
 \draw [BarreStyle=red]  (A-4-1.south west) node[SigneMoins=red] {$-$} to (A-2-3.north east);
 \draw [BarreStyle=red]  (A-5-1.south west) node[SigneMoins=red] {$-$} to (A-3-3.north east);
\end{tikzpicture}
\end{center}
\end{column}
\end{columns}

\end{frame}

\begin{frame}[allowframebreaks]{Linealidad respecto de las filas}

\begin{proposition}
El determinante es lineal respecto de cada fila. Más concretamente,
si \(u,v\in M_{1\times n}(k)\) y \(\lambda\in k\), entonces
\[
\det
\left(
\begin{array}{c}
\vdots\\
\hline
u+v\\
\hline
\vdots
\end{array}
\right)
=
\det
\left(
\begin{array}{c}
\vdots\\
\hline
u\\
\hline
\vdots
\end{array}
\right)
+
\det
\left(
\begin{array}{c}
\vdots\\
\hline
v\\
\hline
\vdots
\end{array}
\right),
\]
y
\[
\det
\left(
\begin{array}{c}
\vdots\\
\hline
\lambda u\\
\hline
\vdots
\end{array}
\right)
=
\lambda
\det
\left(
\begin{array}{c}
\vdots\\
\hline
u\\
\hline
\vdots
\end{array}
\right),
\]
donde las restantes filas son las mismas en todas las matrices.
\end{proposition}

\end{frame}

\begin{frame}[allowframebreaks, t]{Linealidad respecto de la primera fila}

Sean
\[
u=(u_1,\ldots,u_n),\qquad
v=(v_1,\ldots,v_n).
\]
Mantengamos fijas las demás filas de la matriz.

\medskip

Por la definición inductiva del determinante,
\[
\begin{aligned}
\det
\left(
\begin{array}{c}
u+v\\
\hline
\vdots
\end{array}
\right)
&=
\sum_{j=1}^n
(-1)^{1+j}(u_j+v_j)\det(A_{1j})\\
&=
\sum_{j=1}^n
(-1)^{1+j}u_j\det(A_{1j})
+
\sum_{j=1}^n
(-1)^{1+j}v_j\det(A_{1j})\\
&=
\det
\left(
\begin{array}{c}
u\\
\hline
\vdots
\end{array}
\right)
+
\det
\left(
\begin{array}{c}
v\\
\hline
\vdots
\end{array}
\right).
\end{aligned}
\]

Del mismo modo, para \(\lambda\in k\),
\[
\begin{aligned}
\det
\left(
\begin{array}{c}
\lambda u\\
\hline
\vdots
\end{array}
\right)
&=
\sum_{j=1}^n
(-1)^{1+j}(\lambda u_j)\det(A_{1j})\\
&=
\lambda\sum_{j=1}^n
(-1)^{1+j}u_j\det(A_{1j})\\
&=
\lambda\det
\left(
\begin{array}{c}
u\\
\hline
\vdots
\end{array}
\right)
\end{aligned}
\]

Por tanto, el determinante es lineal respecto de la primera fila.

\end{frame}

\begin{frame}[allowframebreaks, t]{Linealidad respecto otra fila}

El resultado es obvio para matrices de orden \(1\). Supongamos que \(n>1\) y que el resultado es cierto para matrices de orden \(n-1\).

\medskip

Supongamos también que la suma de vectores o el producto por un escalar se realiza en otra fila, la fila \(i>1\). 

\medskip

Dado un vector $w=(w_1,\ldots,w_n)$, denotamos 
\[w^j=(w_1,\ldots,w_{j-1},w_{j+1},\ldots,w_n)\]
el vector obtenido suprimiendo su componente \(j\). 

\framebreak

La definición inductiva del determinante y la hipótesis de inducción nos dan:
\begin{multline*}
\det
\left(
\begin{array}{c}
\vdots\\
\hline
u+v\\
\hline
\vdots
\end{array}
\right)
=
\sum_{j=1}^n
(-1)^{1+j}a_{1j}\det
\left(
\begin{array}{c}
\vdots\\
\hline\\[-12pt]
u^j+v^j\\
\hline
\vdots
\end{array}
\right)\\
=
\sum_{j=1}^n
(-1)^{1+j}a_{1j}\left(\det
\left(
\begin{array}{c}
\vdots\\
\hline\\[-12pt]
u^j\\
\hline
\vdots
\end{array}
\right)
+
\det
\left(
\begin{array}{c}
\vdots\\
\hline\\[-12pt]
v^j\\
\hline
\vdots
\end{array}
\right)\right)
\\
=
\sum_{j=1}^n
(-1)^{1+j}a_{1j}\det
\left(
\begin{array}{c}
\vdots\\
\hline\\[-12pt]
u^j\\
\hline
\vdots
\end{array}
\right)
+
\sum_{j=1}^n
(-1)^{1+j}a_{1j}\det
\left(
\begin{array}{c}
\vdots\\
\hline\\[-12pt]
v^j\\
\hline
\vdots
\end{array}
\right)
=
\det
\left(
\begin{array}{c}
\vdots\\
\hline
u\\
\hline
\vdots
\end{array}
\right)
+
\det
\left(
\begin{array}{c}
\vdots\\
\hline
v\\
\hline
\vdots
\end{array}
\right).
\end{multline*}

\framebreak

Del mismo modo, para \(\lambda\in k\),
\[
\begin{aligned}
\det
\left(
\begin{array}{c}
\vdots\\
\hline
\lambda u\\
\hline
\vdots
\end{array}
\right)
&=
\sum_{j=1}^n
(-1)^{1+j}a_{1j}\det
\left(
\begin{array}{c}
\vdots\\
\hline\\[-12pt]
\lambda u^j\\
\hline
\vdots
\end{array}
\right)=
\sum_{j=1}^n
(-1)^{1+j}a_{1j}\lambda\det
\left(
\begin{array}{c}
\vdots\\
\hline\\[-12pt]
u^j\\
\hline
\vdots
\end{array}
\right)\\
&=
\lambda\sum_{j=1}^n
(-1)^{1+j}a_{1j}\det
\left(
\begin{array}{c}
\vdots\\
\hline\\[-12pt]
u^j\\
\hline
\vdots
\end{array}
\right)
=
\lambda\det
\left(
\begin{array}{c}
\vdots\\
\hline
u\\
\hline
\vdots
\end{array}
\right).
\end{aligned}
\]

Por tanto, el determinante es lineal respecto de cualquier fila.

\end{frame}

\begin{frame}[allowframebreaks, t]{Dos filas consecutivas iguales}

\begin{proposition}
Si una matriz cuadrada tiene dos filas consecutivas iguales, entonces
su determinante es \(0\).
\end{proposition}

\begin{proof}
Procedemos por inducción sobre el orden \(n\).

Para \(n=2\),
\[
\det
\begin{pmatrix}
a&b\\
a&b
\end{pmatrix}
=ab-ba=0.
\]

Supongamos ahora \(n>2\) y que el resultado es cierto para matrices
de orden \(n-1\).

Supongamos primero que las filas iguales son las filas \(i\) e
\(i+1\), con \(i>1\). Al desarrollar por la primera fila,
\[
\det(A)
=
\sum_{j=1}^n
(-1)^{1+j}a_{1j}\det(A_{1j}).
\]
En cada \(A_{1j}\), las filas \(i-1\) e \(i\) son iguales y
consecutivas. Por la hipótesis de inducción,
\[
\det(A_{1j})=0
\]
para todo \(j\). Por tanto,
\[
\det(A)=0.
\]

Queda el caso en que las dos primeras filas son iguales. Escribamos
\[
A=
\begin{pmatrix}
a_1&a_2&\cdots&a_n\\
a_1&a_2&\cdots&a_n\\
\multicolumn{4}{c}{\raisebox{1ex}{\(\vdots\)}}
\end{pmatrix}.
\]

Para \(j<k\), denotemos por \(B_{jk}\) la matriz obtenida suprimiendo
las dos primeras filas y las columnas \(j\) y \(k\).

Desarrollamos primero por la primera fila y, en cada una de las
matrices obtenidas, nuevamente por la primera fila.

Para cada par \(j<k\) aparecen exactamente dos términos que contienen
el producto
\[
a_ja_k\det(B_{jk}).
\]

El que se obtiene suprimiendo primero la columna \(j\) y después la
columna \(k\) lleva el signo
\[
(-1)^{1+j}(-1)^k
=
(-1)^{j+k+1},
\]
mientras que el que se obtiene suprimiendo primero la columna \(k\)
y después la columna \(j\) lleva el signo
\[
(-1)^{1+k}(-1)^{1+j}
=
(-1)^{j+k+2}.
\]

Estos signos son opuestos. Por tanto, los términos se cancelan dos a
dos y
\[
\det(A)=0.
\]
\end{proof}

\end{frame}

\begin{frame}[allowframebreaks, t]{Intercambio de dos filas consecutivas}

\begin{proposition}
Si \(A'\) se obtiene de \(A\) intercambiando dos filas consecutivas,
entonces
\[
\det(A')=-\det(A).
\]
\end{proposition}

\begin{proof}
Sean \(u,v\in M_{1\times n}(k)\) las dos filas consecutivas. Como
la matriz
\[
\left(
\begin{array}{c}
\vdots\\
\hline
u+v\\
u+v\\
\hline
\vdots
\end{array}
\right)
\]
tiene dos filas consecutivas iguales, su determinante es \(0\).

\framebreak

Por la linealidad respecto de cada una de esas filas,
\[
\begin{aligned}
0
={}&
\det
\left(
\begin{array}{c}
\vdots\\
\hline
u\\
u\\
\hline
\vdots
\end{array}
\right)
+
\det
\left(
\begin{array}{c}
\vdots\\
\hline
u\\
v\\
\hline
\vdots
\end{array}
\right)+
\det
\left(
\begin{array}{c}
\vdots\\
\hline
v\\
u\\
\hline
\vdots
\end{array}
\right)
+
\det
\left(
\begin{array}{c}
\vdots\\
\hline
v\\
v\\
\hline
\vdots
\end{array}
\right).
\end{aligned}
\]

El primer y el último determinante son \(0\). Por tanto,
\[
\det
\left(
\begin{array}{c}
\vdots\\
\hline
v\\
u\\
\hline
\vdots
\end{array}
\right)
=
-
\det
\left(
\begin{array}{c}
\vdots\\
\hline
u\\
v\\
\hline
\vdots
\end{array}
\right).
\]
\end{proof}

\end{frame}

\begin{frame}{Intercambio de dos filas}

\begin{corollary}
Si \(A'\) se obtiene de \(A\) intercambiando dos filas distintas,
entonces
\[
\det(A')=-\det(A).
\]
\end{corollary}

\vspace{-15pt}

\begin{proof}
Supongamos que se intercambian las filas \(i\) y \(j\), con \(i<j\).

Podemos llevar la fila \(i\) hasta la posición \(j\) mediante
\(j-i\) intercambios de filas consecutivas, y llevar después la fila
que inicialmente ocupaba la posición \(j\) hasta la posición \(i\)
mediante \(j-i-1\) intercambios más.

En total realizamos
\[
(j-i)+(j-i-1)=2(j-i)-1
\]
intercambios de filas consecutivas.

Cada uno cambia el signo del determinante y \(2(j-i)-1\) es impar.
Por tanto,
\[
\det(A')=-\det(A).\qedhere
\]
\end{proof}

\end{frame}

\begin{frame}{Dos filas iguales}

\begin{corollary}
Si una matriz cuadrada $A$ tiene dos filas iguales, entonces
\[
\det(A)=0.
\]
\end{corollary}

\vspace{-15pt}

\begin{proof}
Supongamos que las filas \(i\) y \(j\), con \(i<j\), son iguales. 
Si son consecutivas, el resultado ya está demostrado. Si no, sea
\(A'\) es la matriz que se obtiene de \(A\) intercambiando las filas \(i\) y \(j-1\). Entonces
\[\det(A)=-\det(A'),\]
además \(A'\) tiene dos filas consecutivas iguales, de modo que
\[\det(A')=0.
\]
En consecuencia,
\[\det(A)=0.\qedhere
\]
\end{proof}

\end{frame}

\begin{frame}[allowframebreaks]{Desarrollo por una fila cualquiera}

\begin{proposition}
Sea \(A=(a_{ij})\in M_n(k)\). Para cualquier fila \(i\),
\[
\det(A)
=
\sum_{j=1}^n
(-1)^{i+j}a_{ij}\det(A_{ij}),
\]
donde \(A_{ij}\) es la matriz que se obtiene de \(A\) suprimiendo
la fila \(i\) y la columna \(j\).
\end{proposition}

\begin{proof}
Llevamos la fila \(i\) de \(A\) a la primera posición mediante
\(i-1\) intercambios de filas consecutivas, y denotamos por \(B\)
la matriz obtenida. Entonces
\[
\det(B)=(-1)^{i-1}\det(A).
\]

La primera fila de \(B\) es
\[
(a_{i1},\ldots,a_{in}).
\]
Además, al suprimir de \(B\) su primera fila y la columna \(j\)
obtenemos precisamente \(A_{ij}\).

Por la definición del determinante,
\[
\det(B)
=
\sum_{j=1}^n
(-1)^{1+j}a_{ij}\det(A_{ij}).
\]

Por tanto,
\[
\begin{aligned}
\det(A)
&=
(-1)^{i-1}\det(B)\\
&=
\sum_{j=1}^n
(-1)^{i+j}a_{ij}\det(A_{ij}).
\end{aligned}
\]
\end{proof}

\end{frame}

\begin{frame}[allowframebreaks, t]{Determinantes y operaciones elementales}

\begin{theorem}
Sea \(A\in M_n(k)\).

\begin{enumerate}
\item Si
\(
A\xrightarrow{F_i\leftrightarrow F_j}B,
\)
entonces
\[
\det(B)=-\det(A).
\]

\item Si
\(
A\xrightarrow{\lambda F_i}B\),
\(
\lambda\neq0,
\)
entonces
\[
\det(B)=\lambda\det(A).
\]

\item Si
\(
A\xrightarrow{F_i+\lambda F_j}B\),
\(
i\neq j,
\)
entonces
\[
\det(B)=\det(A).
\]
\end{enumerate}
\end{theorem}

\framebreak 

\begin{proof}
La primera afirmación ya está demostrada.

La segunda se sigue de la linealidad del determinante respecto de
cada fila.

Para la tercera, por linealidad respecto de la fila \(i\),
\[
\det(B)
=
\det(A)+\lambda\det(C),
\]
donde \(C\) se obtiene de \(A\) sustituyendo la fila \(i\) por la
fila \(j\). Como \(C\) tiene dos filas iguales,
\[
\det(C)=0.
\]
Por tanto,
\[
\det(B)=\det(A).
\]
\end{proof}

\end{frame}

\begin{frame}{Diagonal principal}

\begin{definition}
Sea \(A=(a_{ij})\in M_n(k)\). La \alert{diagonal principal} de \(A\)
es la formada por las entradas
\[
a_{11},a_{22},\ldots,a_{nn}.
\]
\end{definition}

\[
A=
\begin{pmatrix}
\alert{a_{11}} & a_{12} & \cdots & a_{1n}\\
a_{21} & \alert{a_{22}} & \cdots & a_{2n}\\
\vdots & \vdots & \ddots & \vdots\\
a_{n1} & a_{n2} & \cdots & \alert{a_{nn}}
\end{pmatrix}
\]

\end{frame}

\begin{frame}{Matrices triangulares superiores}

\begin{definition}
Una matriz \(A=(a_{ij})\in M_n(k)\) se denomina
\alert{triangular superior} si
\[
a_{ij}=0
\qquad\text{si }i>j.
\]

Es decir, tiene la forma
\[
A=
\begin{pmatrix}
a_{11} & a_{12} & \cdots & a_{1n}\\
0      & a_{22} & \cdots & a_{2n}\\
\vdots & \ddots & \ddots & \vdots\\
0      & \cdots & 0      & a_{nn}
\end{pmatrix}.
\]
Dicho de otro modo, todos los elementos por debajo de la diagonal principal son \(0\).
\end{definition}

\end{frame}
\begin{frame}[allowframebreaks]{Determinante de una matriz triangular superior}

\begin{proposition}
Sea
\[
A=
\begin{pmatrix}
a_{11} & a_{12} & \cdots & a_{1n}\\
0      & a_{22} & \cdots & a_{2n}\\
\vdots & \ddots & \ddots & \vdots\\
0      & \cdots & 0      & a_{nn}
\end{pmatrix}
\]
una matriz triangular superior. Entonces
\[
\det(A)=a_{11}a_{22}\cdots a_{nn}.
\]
\end{proposition}

\begin{proof}
Procedemos por inducción sobre \(n\).

Para \(n=1\), el resultado es inmediato.

Supongamos que el resultado es cierto para matrices triangulares
superiores de orden \(n-1\).

Desarrollamos el determinante de \(A\) por la última fila. Como todos
sus elementos salvo el último son \(0\),
\[
\det(A)
=
(-1)^{n+n}a_{nn}\det(A_{nn})
=
a_{nn}\det(A_{nn}).
\]

La matriz \(A_{nn}\) es triangular superior de orden \(n-1\), con
diagonal principal
\[
a_{11},a_{22},\ldots,a_{n-1,n-1}.
\]
Por la hipótesis de inducción,
\[
\det(A_{nn})
=
a_{11}a_{22}\cdots a_{n-1,n-1}.
\]

Por tanto,
\[
\det(A)
=
a_{11}a_{22}\cdots a_{nn}.
\]
\end{proof}

\end{frame}

\begin{frame}{Cálculo del determinante mediante eliminación gaussiana}

Sea \(A\in M_n(k)\). Mediante operaciones elementales de tipos \(1\)
y \(3\), obtenemos una forma escalonada \(E\) de \(A\).

Como \(A\) es cuadrada, \(E\) es triangular superior y, por tanto,
\[
\det(E)=e_{11}\cdots e_{nn}.
\]

Las operaciones de tipo \(3\) no modifican el determinante y cada
intercambio de filas cambia su signo. Por tanto, si se han realizado
\(s\) intercambios,
\[
\det(A)=(-1)^s e_{11}\cdots e_{nn}.
\]

\end{frame}

\begin{frame}{Ejemplo: cálculo del determinante por Gauss}

Sea
\[
A=
\begin{pmatrix}
0&2&1&3\\
1&2&-1&0\\
0&4&1&8\\
0&0&-1&5
\end{pmatrix}.
\]

Mediante eliminación gaussiana,
\[
A
\xrightarrow{F_1\leftrightarrow F_2}
\begin{pmatrix}
1&2&-1&0\\
0&2&1&3\\
0&4&1&8\\
0&0&-1&5
\end{pmatrix}
\xrightarrow{F_3-2F_2}
\begin{pmatrix}
1&2&-1&0\\
0&2&1&3\\
0&0&-1&2\\
0&0&-1&5
\end{pmatrix}
\xrightarrow{F_4-F_3}
\begin{pmatrix}
1&2&-1&0\\
0&2&1&3\\
0&0&-1&2\\
0&0&0&3
\end{pmatrix}.
\]

Se ha realizado un único intercambio de filas. Por tanto,
\[
\det(A)
=
-\,\bigl(1\cdot2\cdot(-1)\cdot3\bigr)
=
6.
\]

\end{frame}

\begin{frame}[allowframebreaks, t]{Determinante e invertibilidad}

\begin{theorem}
Sea \(A\in M_n(k)\). Entonces
\[
A\text{ es invertible}
\quad\Longleftrightarrow\quad
\det(A)\neq0.
\]
\end{theorem}

\begin{proof}
Sea $E$ una forma escalonada de \(A\). Según hemos visto $\det(A)=\pm\det(E)$, con lo cual uno de los dos es nulo si y solo si el otro lo es.

\medskip

Si la fila \(i\)  de $E$ es no
nula, su pivote está en una columna \(j_i\geq i\), pues las posiciones
de los pivotes aumentan estrictamente al bajar de fila.
Por tanto, todas las entradas anteriores al pivote de la fila \(i\)
son nulas y, en particular, $E$ es triangular superior.

\medskip

Si \(A\) no es invertible, entonces \(\operatorname{rg}(A)<n\) y, por tanto $E$ tiene la última fila nula. En consecuencia $e_{nn}=0$ y al ser \(E\) triangular superior, \(\det(E)=e_{11}\cdots e_{nn}=0\).

\medskip

Si \(A\) es invertible, entonces \(\operatorname{rg}(A)=n\) y, por tanto, \(E\) no tiene filas nulas. Como cada pivote de $E$ está más a la derecha que el pivote de la fila anterior, los pivotes de $E$ han  de ser las entradas de la diagonal principal. Luego $e_{ii}\neq0$ para todo \(i=1,\ldots,n\) y $\det(E)\neq 0$.
\end{proof}

\end{frame}

\begin{frame}[allowframebreaks]{Determinante de un producto}

\begin{theorem}
Si \(A,B\in M_n(k)\), entonces
\[
\det(AB)=\det(A)\det(B).
\]
\end{theorem}

\begin{proof}
Supongamos primero que \(A\) no es invertible. Mediante eliminación
gaussiana obtenemos una forma escalonada \(E\) de \(A\) con alguna
fila nula.

Aplicando a \(AB\) las mismas operaciones elementales obtenemos
\(EB\), que también tiene una fila nula. Por tanto,
\[
\det(AB)=0.
\]
Como \(A\) no es invertible, también \(\det(A)=0\), y la fórmula se
cumple.

Supongamos ahora que \(A\) es invertible. Mediante operaciones
elementales por filas llevamos \(A\) hasta \(I_n\). Sea
\(\mu\neq0\) el factor por el que estas operaciones multiplican el
determinante. Entonces
\[
1=\mu\det(A).
\]

Aplicando las mismas operaciones a \(AB\), obtenemos \(B\), y por
tanto
\[
\det(B)=\mu\det(AB).
\]

Como \(\mu=\det(A)^{-1}\), se sigue que
\[
\det(AB)=\det(A)\det(B).
\]
\end{proof}

\end{frame}

\begin{frame}{Matriz traspuesta}

\begin{definition}
Sea
\[
A=(a_{ij})=\begin{pmatrix}
a_{11}&a_{12}&\cdots&a_{1n}\\
a_{21}&a_{22}&\cdots&a_{2n}\\
\vdots&\vdots&&\vdots\\
a_{m1}&a_{m2}&\cdots&a_{mn}
\end{pmatrix}\in M_{m\times n}(k).
\]
La \alert{matriz traspuesta} de \(A\) es la matriz
\[
A^t=(a_{ij}^t)=(a_{ji})=\begin{pmatrix}
a_{11}&a_{21}&\cdots&a_{m1}\\
a_{12}&a_{22}&\cdots&a_{m2}\\
\vdots&\vdots&&\vdots\\
a_{1n}&a_{2n}&\cdots&a_{mn}
\end{pmatrix}\in M_{n\times m}(k).
\]
Es decir,
\[a^t_{ij}=a_{ji}.\]
\end{definition}

\end{frame}

\begin{frame}{Traspuesta de un producto}

\begin{proposition}
Sean
\(
A\in M_{m\times n}(k),
\) y \(
B\in M_{n\times p}(k).
\). 
Entonces
\[
(AB)^t=B^tA^t.
\]
\end{proposition}

\begin{proof}
Para \(1\leq i\leq p\) y \(1\leq j\leq m\),
\[
\begin{aligned}
\bigl((AB)^t\bigr)_{ij}
&=(AB)_{ji}=\sum_{r=1}^n a_{jr}b_{ri}\\
&=\sum_{r=1}^n b_{ri}a_{jr}=\sum_{r=1}^n b^t_{ir}a^t_{rj}=(B^tA^t)_{ij}.
\end{aligned}
\]

Por tanto, como ambos lados son matrices de orden \(p\times m\) y tienen las mismas entradas, coinciden.
\end{proof}

\end{frame}

\begin{frame}{Traspuesta de una matriz invertible}

\begin{corollary}
Sea \(A\in M_n(k)\) una matriz invertible. Entonces \(A^t\) es
invertible y
\[
(A^t)^{-1}=(A^{-1})^t.
\]
\end{corollary}

\begin{proof}
Como \(A\) es invertible,
\[
AA^{-1}=I_n
\qquad\text{y}\qquad
A^{-1}A=I_n.
\]

Trasponiendo y usando la fórmula para la traspuesta de un producto,
obtenemos
\[
(AA^{-1})^t=(A^{-1})^tA^t=I_n
\qquad\text{y}\qquad
(A^{-1}A)^t=A^t(A^{-1})^t=I_n.
\]

Por tanto, \((A^{-1})^t\) es una inversa de \(A^t\). Así, \(A^t\)
es invertible y
\[
(A^t)^{-1}=(A^{-1})^t.\qedhere
\]
\end{proof}

\end{frame}

\begin{frame}{Traspuestas de matrices elementales}

\begin{proposition}
Si \(E\in M_n(k)\) es una matriz elemental, entonces
\[
\det(E^t)=\det(E).
\]
\end{proposition}

\begin{proof}
Para los tres tipos de matrices elementales,
\[
E_{ij}^t=E_{ij},
\qquad
E_i(\lambda)^t=E_i(\lambda),
\qquad
E_{ij}(\lambda)^t=E_{ji}(\lambda).
\]

Las matrices \(E_{ij}(\lambda)\) y \(E_{ji}(\lambda)\) tienen
determinante \(1\), pues se obtienen de \(I_n\) mediante una
operación elemental de tipo \(3\).

Por tanto, en todos los casos,
\[
\det(E^t)=\det(E).
\]
\end{proof}

\end{frame}

\begin{frame}[allowframebreaks]{Determinante de la matriz traspuesta}

\begin{theorem}
Para toda \(A\in M_n(k)\),
\[
\det(A^t)=\det(A).
\]
\end{theorem}

\begin{proof}
Supongamos primero que \(A\) es invertible. Como su forma escalonada
y reducida es \(I_n\), \(A\) es producto de matrices elementales:
\[
A=E_1\cdots E_r.
\]

Entonces
\[
A^t=E_r^t\cdots E_1^t.
\]
Por la fórmula del determinante de un producto y el resultado
anterior,
\[
\begin{aligned}
\det(A^t)
&=\det(E_r^t)\cdots\det(E_1^t)\\
&=\det(E_r)\cdots\det(E_1)\\
&=\det(E_1)\cdots\det(E_r)\\
&=\det(A).
\end{aligned}
\]

Supongamos ahora que \(A\) no es invertible. Entonces \(A^t\)
tampoco es invertible pues si \(A^t\) fuese invertible,
también lo sería
\[
(A^t)^t=A.
\]
En este caso \(\det(A)=0=\det(A^t)\).
\end{proof}

\end{frame}

\begin{frame}{Adjuntos de una matriz}

\begin{definition}
Sea \(A=(a_{ij})\in M_n(k)\). El \alert{adjunto} de la entrada
\(a_{ij}\) es el escalar
\[
(-1)^{i+j}\det(A_{ij}),
\]
donde \(A_{ij}\) se obtiene suprimiendo la fila \(i\) y la columna
\(j\) de \(A\).

La \alert{matriz adjunta} de \(A\) es
\[
\operatorname{adj}(A)=((-1)^{i+j}\det(A_{ij})),
\]
es decir, aquella cuya entrada \((i,j)\) es el adjunto de \(a_{ij}\).
\end{definition}

\end{frame}

\begin{frame}{La matriz adjunta}

\begin{proposition}
Para toda \(A\in M_n(k)\),
\[
A\operatorname{adj}(A)^t=\det(A)I_n.
\]
\end{proposition}

\vspace{-15pt}

\begin{proof}
La entrada \((i,j)\) de \(A\operatorname{adj}(A)^t\) es
\[
\sum_{r=1}^n a_{ir}(-1)^{j+r}\det(A_{jr}).
\]

Si \(i=j\), esta suma es precisamente el desarrollo de
\(\det(A)\) por la fila \(i\). 

Si \(i\neq j\), la fórmula es el desarrollo por la fila \(j\) del determinante de la matriz \(B\) la matriz obtenida de
\(A\) sustituyendo la fila \(j\) por la fila \(i\). Como \(B\) tiene dos filas iguales (la fila \(i\) y la fila \(j\)),
\[
\det(B)=0.
\]
Así, las entradas diagonales de \(A\operatorname{adj}(A)^t\) son
\(\det(A)\) y las restantes son \(0\). 
\end{proof}

\end{frame}

\begin{frame}{Fórmula de la inversa por adjuntos}

\begin{corollary}
Si \(A\in M_n(k)\) es invertible, entonces
\[
A^{-1}
=
\frac{1}{\det(A)}\operatorname{adj}(A)^t.
\]
\end{corollary}


\end{frame}